\chapter{Discussion}

\section{Spatial Non-Uniformity and Its Anatomical Interpretation}

The statistically confirmed non-uniform distribution of scar burden across the AHA segments
is anatomically expected: ischaemic infarction follows coronary territory boundaries, making
a spatially uniform distribution implausible a priori. The value of the formal test is therefore
not surprise but documentation---it establishes that the data are internally consistent with
anatomical priors and provides a quantitative benchmark against which any generative
model must be evaluated. The concentration of burden in the inferior-to-apical
mid-ventricular wall, straddling the RCA--LCX boundary, aligns with the known vulnerability
of this region to multi-territory watershed ischaemia and with electrophysiological mapping
studies that consistently identify the inferoseptal and inferolateral walls as common sites for
re-entrant VT circuits \cite{mahida2017}. The reliable sparing of the basal ring reflects the
lower probability of occlusion at the proximal segments of all three coronary
territories---basal involvement requires either a proximal occlusion or extensive downstream
wavefront propagation, both anatomically less frequent than mid-artery occlusion.

The absence of a significant territory-level burden difference should not be read as evidence
that culprit vessel is geometrically irrelevant. The territory archetype maps demonstrate that
LAD-, RCA-, and LCX-dominant scars occupy distinct spatial footprints that scalar mean
burden cannot distinguish. This dissociation between scalar and geometric information is
precisely the limitation of existing risk stratification approaches highlighted in the
introduction: global metrics such as total scar burden or ejection fraction compress spatial
structure that is mechanistically relevant to arrhythmogenesis \cite{mahida2017}. The
present characterization, by preserving spatial information in the polar-map representation,
is positioned to contribute features that scalar approaches discard.

The high dimensionality of the scar pattern space---six PCA components for 72\% of
variance, PC1 accounting for only 18\%---has a direct mechanistic interpretation: infarct
location, extent, transmurality, and longitudinal distribution vary largely independently across
patients, producing a genuinely high-dimensional population distribution. This stands in
contrast to healthy cardiac anatomy, where statistical shape models typically achieve
compact low-dimensional representations because morphological variability is constrained
by developmental and functional pressures \cite{rodero2021, doste2026}. Post-infarction scar
geometry is not so constrained: the location and extent of ischaemic injury depend on which
vessel occludes, at what level, and for how long---factors that vary independently across
patients. Any generative model that imposes strong low-dimensional structure will
necessarily misrepresent this diversity.

\section{Methodological Choices and Their Consequences}

The need for a multi-candidate apex-detection strategy reflects a broader challenge in
computational cardiac imaging: post-infarction remodelling alters ventricular geometry in
ways that invalidate assumptions that hold in healthy cohorts. A fixed-axis approach,
standard in many shape analysis pipelines, fails in approximately one quarter of this
cohort---a non-negligible error rate that would propagate systematically into all downstream
spatial comparisons. The automated solution developed here trades marginal computational
cost for robustness that is essential at cohort scale.

The comparison with the method of N\'{u}\~{n}ez Garc\'{i}a \cite{nunez2018} reveals two distinct methodological
differences with detectable consequences. The first concerns radial distribution: the $r^{1/3}$
displacement applied by the reference method concentrates coverage toward the basal ring,
while the linear mapping used here distributes it uniformly along the apex-to-base axis,
providing a more direct and geometrically interpretable encoding of longitudinal position. The
second concerns the scar surface sampled: the reference method proximity-labels
endocardial surface points and therefore captures the endocardial face of the scar, whereas
the proposed method projects Core Surface vertices directly and therefore represents the
full transmural extent of the scar mass. This distinction is relevant beyond methodological
comparison: transmurality is one of the geometric features most strongly associated with
arrhythmogenic potential, as deeply penetrating scar creates more complete conduction
barriers and more stable re-entrant pathways \cite{mahida2017, cedilnik2023}. These two
differences are independent of each other and should not be conflated.

For the goals of this work---population-level characterisation of scar geometry and
retrieval-based generation---the proposed method's choices are better suited on three
counts: the linear $s$ coordinate provides a geometrically faithful and monotonic
apex-to-base axis directly usable as a statistical descriptor; projecting Core Surface vertices
preserves the full transmural extent of the scar mass, which is the quantity of interest for
arrhythmia substrate characterisation; and the fully automated apex and axis detection
eliminates the need for manual landmark selection required by the reference
method---three seeds per case---making the pipeline applicable at cohort scale without
operator intervention.

The counter-intuitive degradation of retrieval performance following SVD-based feature
pruning points to a general property of heterogeneous clinical cohorts: discriminative
information can be distributed across many dimensions that individually contribute little
variance. Removing them by a variance threshold discards signal alongside noise. This
finding argues against aggressive dimensionality reduction prior to nearest-neighbour
retrieval in settings where the population is genuinely diverse, and suggests that feature
selection in such contexts should be guided by retrieval performance rather than by variance
decomposition alone.

\section{What Retrieval Quality Reveals About the Cohort}

The leave-one-out retrieval results should be read not only as a performance evaluation but
as a characterisation of cohort structure. The right-skewed distribution of scar burden error
and the long low-overlap tail of segment Dice---with a majority of patients achieving overlap
above 0.5 but a substantial minority falling well below---reflect the heterogeneity documented
in Sections~\ref{sec:scarmorph} and~\ref{sec:terrburden}: approximately 30\% of patients
have single-territory scar, which is well represented in the case base and retrieves
high-quality matches, while the 39\% with two-territory and 13\% with three-territory
involvement occupy sparser regions of descriptor space where few geometrically similar
cases exist among the remaining 104. This is not a failure of the retrieval method but a
faithful reflection of the population: example-based generation is inherently limited by the
diversity of the library, and its performance ceiling is set by cohort size.

The performance gap between deterministic and probabilistic retrieval modes is small in
absolute terms, confirming that the nearest neighbour remains the single best match while
the probabilistic mode adds value by sampling geometric variability among similar
patients---a distinction that becomes practically relevant when the goal is to generate a
distribution of plausible substrates rather than a single representative case. In practice, the
softmax weights are strongly concentrated on the nearest neighbour for most patients: when
the nearest-neighbour distance is substantially smaller than the second-nearest, the top
candidate receives sampling probability close to 1.0, making deterministic and probabilistic
modes effectively equivalent for those cases. Only patients whose top-10 neighbours are
closely spaced in descriptor space receive meaningfully distributed weights across multiple
candidates.

The disparity between segment Dice (mean 0.597, median 0.667) and polar-map Dice
(mean 0.106, median 0.032) warrants explicit interpretation rather than dismissal. Two
patients can share the same set of involved AHA segments---and therefore score perfectly on
segment Dice---while differing substantially in the precise circumferential boundaries, local
density, and transmural depth of scar within those segments. The polar-map Dice penalises
any such sub-segment misalignment. The large gap between the two metrics therefore
quantifies the amount of geometric information that lies below the resolution of the AHA
17-segment model: it is substantial, and it is precisely the information that the polar-map
representation is designed to preserve. The long right tail of polar-map Dice values (75th
percentile 0.189, maximum 0.529) confirms that geometrically close matches do exist within
the cohort for a meaningful subset of patients, and that retrieval performance would improve
with a larger case base.

\section{Transmurality and Fragment Geometry as Independent Descriptors}
\label{sec:transmurality_discussion}

The transmurality analysis reveals a consistent endocardial-dominant scar profile across the
cohort, with mean fragment onset near 30\% of wall thickness and mean transmurality of
approximately 19\%. This is anatomically expected for ischaemic infarction, where occlusion
of epicardial coronary arteries produces a wavefront of necrosis that propagates from the
subendocardium outward, and complete transmural replacement requires sustained
ischaemia. The dense-core / thin-edge transmural gradient confirms that the central scar
mass is more deeply intramural than its periphery---a pattern consistent with the
inward-to-outward wavefront of ischaemic necrosis, where the subendocardial core is the
longest-ischaemic zone and the peripheral border is the last tissue reached by the necrotic
front.

The territory-level transmurality differences deserve attention. RCA-dominant scar
presenting both the highest onset depth and the highest transmurality in this cohort indicates
that inferior infarcts tend to penetrate more deeply through the wall than anterior or lateral
ones. One anatomical explanation is that the inferior wall is thinner than the anterior wall, so
a necrotic wavefront of the same absolute depth represents a higher fractional
transmurality---making inferolateral infarcts geometrically more transmural even without
greater absolute ischaemic depth. This interpretation cannot be confirmed without outcome
data, but it is consistent with the observation that RCA-dominant patients also showed more
circumferential spread across adjacent territories in the archetype maps.

The within-patient Spearman correlation between segment-level extent and transmurality
(median $\rho = 0.57$, IQR [0.36, 0.80], $n = 97$ patients) confirms partial but incomplete
association: in most patients, more-burdened segments tend to penetrate more deeply, but
the relationship is far from deterministic. At the population level, the cross-segment
correlation between mean extent and mean transmurality per AHA segment is near-zero
($\rho = 0.14$, $p = 0.59$), indicating that segments with high average burden across the
cohort are not systematically more transmural. A minority of patients (4/97) show negative
within-patient correlation, consistent with co-occurrence of extensive subendocardial scar in
one territory and focal but deeply transmural scar in another. The two descriptors are
therefore not interchangeable, and their joint inclusion in the retrieval feature vector is
justified by the retrieval performance results, which showed that pruning features with low
individual variance contributions consistently degraded rather than improved matching
quality.

Fragment shape descriptors capture a further dimension of scar geometry not encoded by
either territory assignment or transmurality. The predominance of irregular, non-compact
lesions (compactness below 0.4 in most patients) reflects the morphological complexity that
makes direct 3D-to-3D scar comparison across patients impractical without a common
coordinate frame. The association between higher total burden and wider circumferential
spread---with the mid-cavity ring showing the greatest tangential extent at the population
level---indicates that large scars tend to wrap around the ventricular wall rather than penetrate
deeply through it at any single location, consistent with the subendocardial wavefront model
of ischaemic injury. This characterises individual scar morphology along the circumferential
axis and is distinct from the inter-segment correlation structure, which concerns the
longitudinal axis: within a territory, segments at different ring levels co-vary more strongly
than circumferentially adjacent segments across territory boundaries, reflecting contiguous
propagation along the apex-to-base direction rather than lateral spread. The two findings
together suggest that scar geometry is organised differently along each spatial axis---circumferential
extent scales with total burden, longitudinal extent is constrained by vascular territory, and
transmural depth remains relatively limited in most patients.

The apex-to-base gradient, with near-zero population mean but extreme individual variability
(slopes ranging from $-0.82$ to $+0.20$), illustrates a recurring theme of this analysis:
population-level averages conceal individual-level structure that is precisely the information a
generative model must reproduce. A generative model that samples each descriptor
independently from its marginal distribution---generating total burden, territory, and gradient
direction without constraint---would produce cases that violate the joint dependencies
observed in real patients: for example, a case assigned high LAD burden but a
base-dominant gradient, or high RCA burden with an anteroseptal spatial footprint. These
combinations do not occur in clinical data because coronary territory, scar location, and
longitudinal distribution are not independent.

\section{Limitations}

Several limitations of this work must be acknowledged.

The morphological diversity documented in the cohort---spanning all coronary territories,
longitudinal levels, and transmural profiles---is not incidental noise but reflects the genuine
variability of post-infarction remodelling across patients. It is therefore treated as a
characterisation of the problem rather than a limitation of the dataset: any geometric
representation or generative framework must accommodate this range rather than assume a
more homogeneous distribution.

The cohort of 105 patients, while sufficient for the population-level characterisation
conducted here, constrains the retrieval-based generation framework in two ways: rare scar
phenotypes are underrepresented in the case base, and the leave-one-out evaluation is
optimistic relative to prospective deployment where query patients may fall outside the
observed distribution entirely. The subgroup analyses involving age ($n = 48$) and sex
($n = 48$, with severe male skew) are explicitly exploratory and should not be interpreted as
establishing the absence of demographic associations; they establish only that the present
cohort is underpowered to detect them.

The polar-map representation, while enabling population-level comparison, introduces
geometric distortions inherent to any cylindrical projection. The bull's-eye representation
compresses apical regions toward the centre of the disk, potentially underrepresenting apical
scar extent in area-based metrics---though the vertex-based primary representation is less
sensitive to this distortion. The linear longitudinal mapping produces uniform bin density
along the apex-to-base axis but places less weight near the basal ring relative to a cube-root
mapping, a trade-off documented in the comparison with the N\'{u}\~{n}ez et al.\ reference
method.

The transmurality computation relies on nearest-neighbour distances to the endocardial and
epicardial surfaces, which treats the wall as locally planar. In regions of high curvature---particularly
the apex and the lateral wall---this approximation introduces error that is difficult to quantify
without ground-truth wall-thickness measurements. The discrete layer meshes exported by
ADAS 3D (Layer\_10 through Layer\_90) provided insufficient resolution for the primary
transmurality analysis and were superseded by the continuous vertex-depth approach, but
the latter is itself dependent on the quality of the LV shell segmentation.

Finally, the absence of paired electrophysiological or outcome data in the geometrically valid
cohort means that none of the geometric descriptors developed here can be validated against
arrhythmia endpoints within this work. The characterisation is purely morphological: whether
the spatial patterns identified---the inferior-to-apical hotspot, the RCA transmurality excess,
the fragment shape associations---translate into differential arrhythmic risk in this specific
cohort remains an open question that the DEVELOP registry's prospective follow-up data are
positioned to answer.
